\section{Conclusion}
\label{sec:conclusion}

We set out to measure how much geomagnetic storms speed up the decay of objects in low Earth
orbit, using one year of public orbit data and treating each storm as a natural experiment. The
estimate is causal under the assumption that no time-varying confounding remains after we adjust
for F10.7 and season. Storms speed up decay by 14.1~m/day over days 0--2 after onset
(0.39~m/day per nT of ap), about 28\% faster than quiet-time decay, and the estimate holds
without the two largest storms and without the outlier filter. Over days 0--2, NRLMSIS~2.1 shows
little pooled error in ordinary storms but under-predicts decay at peak activity, and in the
largest storms its errors change sign from one storm to the next. Vulnerability scales with drag
sensitivity times an altitude factor that falls about threefold from 620 to 750~km. In an
exploratory analysis, Starlink orbit changes are consistent with an approximately complete
offset of the predicted storm loss (1.17, 95\% CI 0.97--1.36), almost all during the storm. The
storm effect is backed by placebo dates, a pre-trend test, a high-altitude control group and a
simulator with a known answer. The vulnerability map is backed by seed stability and synthetic
recovery. The offset ratio is backed only by subset checks, because its placebo is imprecise.

\section*{Author Contributions}
P.~Shiv built the orbit cohort, TLE cleaning and precise-orbit validation. S.~Akash built the
space-weather data and storm catalogue, the physics baseline and outcomes, the physics simulator,
the object catalogue, the active-satellite table, the model-based offset ratio and the
extreme-storm check. S.~Anil designed the analysis-table schema and causal graph and built the
synthetic data, the stacked event study and its checks, and the overlap correction of the offset
model. S.~Holla built the double machine learning, refutations, storm-class analysis and
vulnerability maps. All authors wrote the paper.

\section*{Code and Data Availability}
All inputs are public: TLEs from Space-Track, space weather from NASA OMNIWeb, object data from
CelesTrak and ESA DISCOSweb, and precise orbits from the GRACE-FO and Swarm missions. The code
(Python) and the processed analysis table are provided in the project submission folder. The
notebook \texttt{ICARUS\_final.ipynb} reproduces every headline number and figure from the
processed table; the raw-data download and table build are listed step by step in the
repository README.
